\documentclass[10pt,a4paper]{amsart}
\usepackage[margin=19mm]{geometry}
\usepackage{amsmath,amssymb,mathtools}
\usepackage[scr=rsfs,cal=euler]{mathalfa}
\usepackage{xfrac}
\usepackage[unicode,hidelinks,bookmarks=false]{hyperref}
\newcommand{\ie}{i.e.\ }
\newcommand{\cf}{cf.\ }
\newcommand{\resp}{respectively\ }
\newcommand{\Subsection}{\subsection}
\providecommand{\nobreakdash}{\nobreak\hbox{-}}
\setlength{\emergencystretch}{2em}
\multlinegap0pt
\usepackage{bm}
\usepackage{bbm}
\usepackage{dsfont}
\usepackage{xspace}
\usepackage{cancel}
\usepackage{mathtools}
\usepackage{amsthm}
\makeatletter
\@ifundefined{theorem}{\newtheorem{theorem}{Theorem}}{}
\makeatother
\newcommand{\K}{\mathbb{K}}
\title[On Buzzard's Theoren]{On Buzzard's Theoren}
\author{Zbigniew Jelonek, Gustavo Menani, Maria Michalska}
\date{Source: arXiv:2608.17166v1 (CC BY 4.0)}
\begin{document}
\maketitle

\noindent\emph{Source and licence.} On Buzzard's Theoren. Version arXiv:2608.17166v1 (CC BY 4.0), distributed under the Creative Commons Attribution 4.0 International licence, as recorded in the arXiv licence metadata for this exact version and shown on the abstract page https://arxiv.org/abs/2608.17166v1. Full rights notice: licenses/arxiv-2608.17166-CC-BY-4.0.txt.\par\smallskip

\noindent\textit{Editorial scope.} Counted unit: the Theorem whose source label is \texttt{thm:pointautomorphism} in the article (source0.tex lines 167--172, source label \texttt{thm:pointautomorphism}), delivered together with the article's own complete original proof of it (source0.tex lines 175--197, one \texttt{proof} environment of 23 lines). No mathematics is written, repaired or re-derived: statement and proof are the article's own text line for line. Required context retained uncounted: glowne (lines 120--124). Every definition, display and label of those lines is retained unchanged. Omitted from this package and not redistributed: 1--119, 125--166, 173--174, 198--353. Nothing inside a retained excerpt is omitted or altered: every retained body is delivered exactly as the article's lines are written, so no omission receipt of any kind accompanies this package. Citations: the retained excerpts carry no citation command at all, so no bibliography entry of the article is transcribed and no reference is redistributed. Reference adaptation: none was needed and none was made; every reference inside the delivered unit resolves inside it, so no unresolved reference or citation is produced. Printed slips kept literal: no printed word, symbol, hypothesis or number of the counted statement or of its proof was altered. Layout and class: the article's own class is \texttt{amsart} with packages including amssymb, mathrsfs, hyperref, verbatim, graphicx, subfigure, psfrag, enumitem, xcolor; the wrapper is the lane's pinned \texttt{amsart} editorial wrapper at 10pt on A4, which restates the theorem-like environments the retained text needs and adds the article's own macro definitions that the retained text needs (\texttt{\textbackslash K}), with extra wrapper packages (bm, bbm, dsfont, xspace, cancel, mathtools). The delivered numbering is the unit's own. Assets: no retained span contains a figure environment, a graphics-inclusion command, a PDF-page inclusion, a table or a TikZ picture; no image file is copied into this package, the package declares no asset dependency and no image file is redistributed with it. Licence: version arXiv:2608.17166v1 of this article is distributed under the Creative Commons Attribution 4.0 International licence, as recorded in the arXiv licence metadata for this exact version and shown on the abstract page https://arxiv.org/abs/2608.17166v1. A full rights notice is delivered at licenses/arxiv-2608.17166-CC-BY-4.0.txt. Characters: every retained excerpt is pure ASCII, so the delivered engine, XeLaTeX with its default Latin Modern text font, reports no missing character.
\begin{theorem}\label{glowne}
Let $\K$ be an infinite field. Let  $n\ge 2$ and $a_1,\ldots, a_k; $ $b_1,\ldots , b_k$ be two sequences of different points from $\K^n.$ Let $L_1,\ldots, L_k\in SL(n,\K)$ be linear isomorphisms with 
jacobian one. 
Then there is a polynomial automorphism $\Phi: \K^n\to \K^n$ with jacobian one, such that $\Phi(a_i)=b_i$ and $d_{a_i} \Phi= L_i$ for $i=1,\ldots, k.$
\end{theorem}

\begin{theorem}\label{thm:pointautomorphism}
Let $\K$ be an infinite field. Let $n\ge 2$. 
For any two sets   $\{a^1,\ldots,a^k\}$ and $\{b^1,\ldots,b^k\}$  of $k$ distinct points in~$\K^n$ there exists a polynomial automorphism $F:\K^n \to \K^n$ of degree at most $k^2$ and jacobian determinant $1$ such that
$$ F(a^i) = b^i $$
for each $i=1,\dots,k$.
\end{theorem}

\begin{proof}

It is easy to see that there exist linear isomorphisms $A,B\in SL_n(\K)$ such that the points $A(a^1),\dots, A(a^k)$  are in a general position such that no two points have the same coordinate, i.e.  $a_i^j\neq a_i^l$ for all $i\leq n$, $j, l\leq k$ and $j\neq l$, and the points $B(b^1),\dots,B(b^k)$ are such that no two first coordinates are the same. If there exists an automorphism $F$ satisfying the claim in these new coordinates, then $B^{-1}\circ F\circ A$ satisfies the claim in old coordinates. Thus without loss of generality we may assume that the sets  $\{a^1,\ldots,a^k\}$ and $\{b^1,\ldots,b^k\}$ are in a general position.

By Corollary \ref{glowne} for $1\leq j \leq n-1$ exist polynomials $p_j\in \K[t]$ of degree at most~$k$  such that 
$$p_j(a^{i}_{j+1})=b^{i}_{j}-a^{i}_{j}$$
for every $i=1,\dots, k$.
Let
$$P(x_1,\dots, x_n):=(x_1+p_{1}(x_{2}),\dots,\ x_{n-1}+p_{n-1}(x_{n}),\ x_n).$$
Then $P$ is a tame automorphism and $P(a^i)=(b^{i}_{1},\dots, b^{i}_{n-1}, a^{i}_{n}).$


Let $p_n\in\K[t]$ be such that 
$$p_n(b^{i}_{1})=b^{i}_{n}-a^{i}_{n}$$
for every $i=1,\dots,k$. 
Take
$$H(x_1,\dots,x_n):=(x_1,\dots,x_{n-1},\ x_n+p_n(x_1)).$$
Then  $H$ is a tame polynomial automorphisms of $\K^n.$ Let $$F=  H\circ P.$$ 
We have
$$F(a^i)=  (H\circ P)(a^i)=H(b^{i}_{1},\dots, b^{i}_{n-1}, a^{i}_{n})=(b^{i}_{1},\dots, b^{i}_{n})=b_i $$
%Thus $H\circ P(a_i)=b_i$
 for $i=1,\dots, k.$ Since $\deg P, \deg H\le k$ we have $\deg F\le k^2.$
\end{proof}

\end{document}
